Un model simplificat de distribució de càrregues elèctriques a l'interior d'un núvol es pot aproximar a dues càrregues puntuals situades a diferents altures. La figura mostra aquesta distribució aproximada, que consta d'una càrrega de $+40\ \mathrm{C}$ situada a 10~km d'altura i una càrrega de $-30\ \mathrm{C}$ situada a 4~km d'altura.

\figura[0.35]{fig1.png}

\begin{apartat}{1}
Calculeu el vector camp elèctric que crea el núvol en el punt P que s'indica a la figura.
\end{apartat}

\begin{apartat}{1}
Calculeu l'energia potencial electroestàtica emmagatzemada en el núvol.
\end{apartat}

\blocetiquetat{Dada:}{$k = \dfrac{1}{4\pi\varepsilon_0} = 9,0\times 10^{9}\ \mathrm{N\,m^2\,C^{-2}}$.}
